\documentclass[11pt,a4paper]{article}
\usepackage[T1]{fontenc}
\usepackage{amsmath,amssymb}
\usepackage[margin=1in]{geometry}
\usepackage{hyperref}
\usepackage{xurl}

\title{OpenCV Camera Calibration Distortion Model}
\author{Quest Shooting Demo}
\date{}

\newcommand{\clamp}{\operatorname{clamp}}

\begin{document}
\maketitle

\section{Scope}
This document describes the OpenCV five-coefficient camera-distortion model
used by the project's screen-space shader. It modifies the rendered virtual
scene; it does not access the Quest camera feed or model Quest passthrough.
It is not a headset display-lens correction. All coefficients default to zero
because no calibration measurements have been supplied.

\section{What: the optical model and its equations}
\subsection{The ideal pinhole camera}
In an ideal pinhole camera, a scene point with camera-space coordinates
$(X,Y,Z)$ projects to the normalized image plane at
\begin{equation}
  x = X/Z, \qquad y = Y/Z.
\end{equation}
The coordinates $x$ and $y$ are dimensionless. With focal lengths $f_x,f_y$
in pixels and principal point $(c_x,c_y)$, the ideal pixel position is
\begin{equation}
  u = f_x x + c_x, \qquad v = f_y y + c_y.
\end{equation}
The principal point is the image location associated with the camera's optical
axis in this model; the focal lengths scale normalized coordinates into pixels.
Straight lines in the scene that do not pass through the camera center project
as straight lines in this ideal model.

\subsection{Five-coefficient distortion}
Real lenses can map rays to image positions differently from an ideal
pinhole. Let $(x,y)$ be the ideal normalized coordinates and
$r^2=x^2+y^2$ the squared distance from the optical axis. OpenCV's
five-coefficient model maps the ideal point to the distorted point $(x_d,y_d)$:
\begin{align}
  R(r) &= 1+k_1r^2+k_2r^4+k_3r^6,\\
  x_d &= xR(r)+2p_1xy+p_2(r^2+2x^2),\\
  y_d &= yR(r)+p_1(r^2+2y^2)+2p_2xy.
\end{align}
The distorted pixel location is $u_d=f_xx_d+c_x$,
$v_d=f_yy_d+c_y$. The coefficients are dimensionless under this normalized
coordinate convention. They describe a mapping in the image plane; they are
not lens radii, focal lengths, or physical displacements in millimeters.
All-zero coefficients give $R=1$ and $(x_d,y_d)=(x,y)$.

\section{How: what each coefficient does}
\subsection{Radial distortion: $k_1$, $k_2$, and $k_3$}
The radial terms multiply both image coordinates by the same scale $R(r)$.
For radial distortion alone, when $R(r)>0$, the point stays on the ray from
the image center, but its distance changes from $r$ to $r_d=rR(r)$. Hence its
radial displacement is
\begin{equation}
  r_d-r=k_1r^3+k_2r^5+k_3r^7.
\end{equation}
\begin{itemize}
  \item $k_1$ controls the lowest-order change of magnification with distance
  from the center. It often has the largest visible effect near the center,
  although the relative effect depends on all coefficients and on $r$.
  \item $k_2$ adds a higher-order adjustment, proportional to $r^5$ in
  radial displacement. It changes how the curvature evolves farther from the
  center.
  \item $k_3$ adds another higher-order adjustment, proportional to $r^7$.
  Its effect grows faster with radius than the $k_1$ and $k_2$ terms, for
  coefficients of comparable magnitude.
\end{itemize}
For example, with $k_1=0.1$, $k_2=k_3=p_1=p_2=0$, a point at
$(x,y)=(0.5,0)$ has $R=1.025$ and moves to $(0.5125,0)$. A positive
radial displacement produces a pincushion-like mapping; a negative one
produces a barrel-like mapping. These descriptions depend on the net radial
scale over the field of view, not on the sign of one coefficient alone.

\subsection{Tangential distortion: $p_1$ and $p_2$}
The tangential, or decentering, contribution is the difference
\begin{align}
  \Delta x_t &= 2p_1xy+p_2(r^2+2x^2),\\
  \Delta y_t &= p_1(r^2+2y^2)+2p_2xy.
\end{align}
\begin{itemize}
  \item $p_1$ controls one direction of decentering. It contributes
  $2p_1xy$ horizontally and $p_1(r^2+2y^2)$ vertically.
  \item $p_2$ controls the other direction. It contributes
  $p_2(r^2+2x^2)$ horizontally and $2p_2xy$ vertically.
\end{itemize}
These terms can make the displacement asymmetric about the image center.
Neither coefficient acts only on one image axis. For example, with
$p_1=0.1$ and all other coefficients zero, $(0,0.5)$ maps to $(0,0.575)$.
The two components can combine with each other and with radial distortion.

\section{Why: the optical interpretation}
Lens distortion is a \emph{geometric} aberration: image magnification varies
across the field, shifting image points away from their ideal pinhole
positions. A rotationally symmetric lens system can be approximated by a
radial scale that depends on distance from the optical axis. The powers
$r^2$, $r^4$, and $r^6$ provide successive terms for fitting how that scale
changes across the image. Decentering or misalignment of lens elements can
break radial symmetry; the two tangential terms provide a simple model for
part of that asymmetric displacement.

These are mathematical approximations fitted to observed image points, not a
ray trace or a unique physical reconstruction of the lens. During calibration,
known pattern points are matched to measured image locations and parameters are
estimated to reduce reprojection error. Different lens designs, focus settings,
and image geometry can require different coefficients or a more complex model.
Good calibration images and residual checks are needed to judge whether this
five-coefficient model fits a particular camera. Blur, chromatic aberration,
noise, and Quest passthrough reconstruction are not explained by these
geometric terms.

The panel displays radial terms first: $k_1,k_2,k_3,p_1,p_2$. When supplying
an OpenCV five-value distortion vector, its standard ordering is instead
$(k_1,k_2,p_1,p_2,k_3)$.

\section{Inverse mapping for the screen shader}
The forward calibration map sends a source point to its distorted image
position. A screen shader instead starts with an output pixel and must find
the corresponding source pixel. Define
\begin{equation}
  \mathbf{d}(x,y)=
  \begin{bmatrix}x_d(x,y)\\y_d(x,y)\end{bmatrix},
  \qquad
  \mathbf{q}=
  \begin{bmatrix}x_d^{\mathrm{target}}\\y_d^{\mathrm{target}}\end{bmatrix}.
\end{equation}
The desired source coordinate satisfies
\begin{equation}
  \mathbf{d}(\mathbf{x})=\mathbf{q}.
\end{equation}
The shader solves this with eight Newton iterations, initialized at
$\mathbf{x}_0=\mathbf{q}$:
\begin{align}
  \mathbf{r}_n &= \mathbf{d}(\mathbf{x}_n)-\mathbf{q},\\
  \mathbf{x}_{n+1} &=
  \mathbf{x}_n-\mathbf{J}(\mathbf{x}_n)^{-1}\mathbf{r}_n,
\end{align}
where $\mathbf{J}$ is the Jacobian of the forward map. If its determinant is
too close to zero, that iteration is skipped. Newton's method is an
approximation: eight iterations do not guarantee a solution, especially for
strong distortion or a mapping that is not one-to-one.

For completeness, write
\begin{align}
  s &= 1+k_1r^2+k_2r^4+k_3r^6,\\
  g &= k_1+2k_2r^2+3k_3r^4.
\end{align}
The Jacobian used by the shader is
\begin{equation}
\mathbf{J}=
\begin{bmatrix}
s+2x^2g+2p_1y+6p_2x &
2xyg+2p_1x+2p_2y\\
2xyg+2p_1x+2p_2y &
s+2y^2g+6p_1y+2p_2x
\end{bmatrix}.
\end{equation}
The shader converts screen UV to normalized camera coordinates using the
projected view center and per-eye projection focal scales. It converts the
recovered undistorted coordinate back to source UV and clamps the texture
lookup to the image bounds.
The per-eye projection parameters here are taken from the virtual camera;
they are not measured Quest passthrough-camera intrinsics. The post-process
warps the rendered image, not the physical objects in the scene.

\section{Panel controls}
The in-world panel exposes all five coefficients. Its ranges and increments
are editing conveniences, not limits specified by OpenCV:
\begin{center}
\begin{tabular}{c|c|c|c}
Coefficient & Range & Fine step & Coarse step\\
\hline
$k_1$ & $[-1,1]$ & $0.01$ & $0.1$\\
$k_2$ & $[-1,1]$ & $0.01$ & $0.1$\\
$k_3$ & $[-1,1]$ & $0.01$ & $0.1$\\
$p_1$ & $[-1,1]$ & $0.01$ & $0.1$\\
$p_2$ & $[-1,1]$ & $0.01$ & $0.1$
\end{tabular}
\end{center}
All coefficient rows show the same controls from left to right:
\texttt{-0.1}, \texttt{-0.01}, the current value,
\texttt{+0.01}, \texttt{+0.1}. The outer buttons apply the coarse step and the
buttons beside the value apply the fine step. These ranges and steps are UI
choices, not limits specified by OpenCV. Use coefficients estimated from a
calibration dataset for the intended camera and image geometry. Arbitrarily
chosen coefficients only create a visual effect and may make the inverse
mapping poorly conditioned or visibly distorted.

\section{Reference}
The equations follow OpenCV's
\href{https://github.com/opencv/opencv/blob/4.x/doc/py_tutorials/py_calib3d/py_calibration/py_calibration.markdown}
{camera calibration tutorial}.
The extended model and camera-coordinate conventions are described in the
\href{https://docs.opencv.org/4.x/d9/d0c/group__calib3d.html}
{OpenCV camera calibration and 3D reconstruction reference}.

\end{document}
